% !TeX root = ../navier-stokes-manuscript.tex
\subsection{The estimate that must survive the endpoint}\label{sub:BodyCompatibleIntersection}

At \((N,M)=(0,1)\), the velocity is controlled one spatial step above \(H^1\)
and its time response one step below.  Together they can reach one time slice.
This lowest rectangle asks for
\begin{equation}\label{eq:BodyFirstWholeTerminalRectangle}
  u\in L^2(I_*;H^2),
  \qquad
  u_t\in L^2(I_*;L^2).
\end{equation}

Which part of the equation could produce these two controls?  Set
\begin{equation}\label{eq:BodyEnstrophyDissipationQuantities}
\begin{aligned}
  Z(t)&:=\norm{\Lambda u(t)}_2^2=\norm{\omega(t)}_2^2,
  &
  Y(t)&:=\norm{\Lambda^2u(t)}_2^2=\norm{\nabla\omega(t)}_2^2,
  \\
  D(t)&:=\norm{\Lambda^{3/2}u(t)}_2^2.
\end{aligned}
\end{equation}
\Needspace{10\baselineskip}
The enstrophy balance \eqref{eq:BodyEnstrophyBalance} contains the stretching
term that must be controlled.  Since \(\nabla u\) is recovered from \(\omega\)
by Riesz transforms, H\"older followed by the Sobolev embeddings produces
\begin{align}
  \left|\int_{\R^3}\omega\cdot S\omega\,dx\right|
  &\leq
  \norm{\omega}_2\norm{\omega}_3\norm{S}_6
  \notag\\
  &\leq C Z^{1/2}D^{1/2}Y^{1/2}.
  \label{eq:BodyFirstRectangleStretchingEstimate}
\end{align}
Young's inequality then turns the enstrophy identity into
\begin{equation}\label{eq:BodyEnstrophyFromCriticalDissipation}
  Z'(t)+\nu Y(t)
  \leq C\nu^{-1}D(t)Z(t).
\end{equation}
The coefficient multiplying \(Z\) has now been isolated.  If
\begin{equation}\label{eq:BodyCriticalDissipationInterface}
  D_*:=\int_0^{T_*}D(t)\,dt
  =\int_0^{T_*}\norm{\Lambda^{3/2}u(t)}_2^2\,dt
  <\infty,
\end{equation}
then Gronwall gives
\begin{equation}\label{eq:BodyEnstrophyAndHtwoFromDissipation}
  \sup_{0<t<T_*}Z(t)+\int_0^{T_*}Y(t)\,dt<\infty.
\end{equation}
The first term supplies \(u\in L^\infty(I_*;H^1)\), and the second, together
with the kinetic-energy row, supplies \(u\in L^2(I_*;H^2)\).  Applying the
Leray projection to the momentum equation gives
\[
  u_t=\nu\Delta u-\Leray\bigl((u\cdot\nabla)u\bigr),
\]
while
\begin{equation}\label{eq:BodyNonlinearityFromCriticalDissipation}
  \norm{\Leray((u\cdot\nabla)u)}_2^2
  \leq C\norm{u}_6^2\norm{\nabla u}_3^2
  \leq C ZD.
\end{equation}
Thus \(D_*<\infty\) also gives \(u_t\in L^2(I_*;L^2)\), which completes
\eqref{eq:BodyFirstWholeTerminalRectangle}.  In the other direction, the
cutoff-uniform \(L^2H^2\) row and monotone convergence give
\(u\in L^2(I_*;H^2)\), hence \(D_*<\infty\).  Using the energy identity in
\Cref{prop:ClassicalKineticEnergyIdentity} for the lower
Sobolev pieces, we have reached the exact equivalence
\begin{equation}\label{eq:BodyFirstRectangleCriticalDissipationEquivalence}
  D_*<\infty
  \quad\Longleftrightarrow\quad
  \sup_{0<T<T_*}\mathfrak R_{0,1}(u;T)<\infty.
\end{equation}

The critical-height balance gives the same interface in its original
coordinates.  Because \(D=2E_{3/2}\), integrating
\eqref{eq:BodyFirstCriticalRow} yields
\begin{equation}\label{eq:BodyCriticalPrefixIdentity}
  H(t)+\nu\int_0^tD(s)\,ds
  =H(0)+\int_0^t\mathcal P_{1/2}(s)\,ds.
\end{equation}
Since \(H\geq0\), finiteness of the signed production prefix makes \(D_*\)
finite.  Conversely, \(D_*<\infty\) gives the first rectangle and hence a
bounded \(H\); the same equality then makes the prefix finite.  Hence
\begin{equation}\label{eq:BodyFirstRectangleSignedPrefixEquivalence}
  \sup_{0<t<T_*}\int_0^t\mathcal P_{1/2}(s)\,ds<\infty
  \quad\Longleftrightarrow\quad
  D_*<\infty
  \quad\Longleftrightarrow\quad
  \sup_{0<T<T_*}\mathfrak R_{0,1}(u;T)<\infty.
\end{equation}
The signed-prefix statement is another exact reading of the first rectangle;
it is not a smaller estimate that produces the rectangle for free.

There is one genuine class of data for which this interface closes.  Its
boundary sits inside the critical balance itself.  Put
\[
  a(t):=\norm{\Lambda^{1/2}u(t)}_2,
  \qquad
  d(t):=\norm{\Lambda^{3/2}u(t)}_2^2.
\]
For
\(
  \mathcal B(u,u):=\Leray((u\cdot\nabla)u)
\), the dual Sobolev embeddings give
\begin{align}
  \norm{\mathcal B(u,u)}_{\dot H^{-1/2}}
  &\leq C\norm{u\cdot\nabla u}_{3/2}
  \notag\\
  &\leq C\norm{u}_3\norm{\nabla u}_3
  \leq C a(t)d(t)^{1/2}.
  \label{eq:BodyCriticalNonlinearDualEstimate}
\end{align}
The critical production pairs this source with one further half derivative of
the velocity, so
\begin{equation}\label{eq:BodyCriticalProductionAbsoluteEstimate}
  |\mathcal P_{1/2}(t)|
  \leq C a(t)d(t).
\end{equation}
Inserted into \eqref{eq:BodyFirstCriticalRow}, this becomes
\begin{equation}\label{eq:BodySmallCriticalDataAbsorption}
  \frac12\frac d{dt}a(t)^2
+\bigl(\nu-Ca(t)\bigr)d(t)
\leq0.
\end{equation}
If \(a(0)<\nu/C\), continuity prevents \(a\) from reaching the viscous
threshold: until a first crossing, the coefficient is positive and \(a\) is
nonincreasing, which rules out that crossing.  Consequently
\begin{equation}\label{eq:BodySmallCriticalDataDissipation}
  \frac12a(t)^2
+\bigl(\nu-Ca(0)\bigr)\int_0^t d(s)\,ds
\leq\frac12a(0)^2.
\end{equation}
This pays \(D_*\) from the datum in the small-critical-data regime.  For an
arbitrary Schwartz datum, \(a(0)\) need not lie below the threshold.  The source
and the dissipation have the same critical scaling, so shrinking the picture
does not improve the coefficient in
\eqref{eq:BodySmallCriticalDataAbsorption}.

Energy can still carry the same history into that regime if the classical
branch survives long enough.  The resulting time does not predict when a
singularity occurs.  It is a horizon beyond which a finite maximal time cannot
sit.

\begin{corollary}[Small-data capture horizon]
\label[corollary]{cor:BodySmallDataCaptureHorizon}
Let \(C\) be the constant in
\eqref{eq:BodySmallCriticalDataAbsorption}.  If a nonzero datum
\(u_0\in\Ssigma\) generates a finite maximal time, then
\begin{equation}\label{eq:BodySmallDataCaptureHorizon}
  T_*
  \leq
  \frac{C^4\norm{u_0}_2^4}{2\nu^5}.
\end{equation}
\end{corollary}

\begin{proof}
Fix \(0<T<T_*\).  The kinetic-energy identity in
\Cref{prop:ClassicalKineticEnergyIdentity} gives
\[
  \int_0^T Z(t)\,dt
  \leq\frac{\norm{u_0}_2^2}{2\nu},
\]
so at some \(t_0\in(0,T)\),
\[
  Z(t_0)\leq\frac{\norm{u_0}_2^2}{2\nu T}.
\]
Interpolation between the energy and enstrophy heights then gives
\[
  a(t_0)^2
  \leq\norm{u(t_0)}_2\norm{\Lambda u(t_0)}_2,
  \qquad
  a(t_0)
  \leq\frac{\norm{u_0}_2}{(2\nu T)^{1/4}}.
\]
If \(T>C^4\norm{u_0}_2^4/(2\nu^5)\), this places the shifted history below
the threshold \(\nu/C\).  The same first-crossing argument, now begun at
\(t_0\), gives
\[
  \frac12a(t)^2
  +\bigl(\nu-Ca(t_0)\bigr)\int_{t_0}^t d(s)\,ds
  \leq\frac12a(t_0)^2,
  \qquad t_0\leq t<T_*.
\]
This pays the remaining critical dissipation.  The portion before \(t_0\) is
finite by classical smoothness, so \(D_*<\infty\).  By
\eqref{eq:BodyEnstrophyAndHtwoFromDissipation}, the same history is then
bounded in \(H^1\), hence in \(L^3\), all the way to the alleged last time.
That contradicts the Escauriaza--Seregin--\v{S}ver\'ak endpoint criterion
\eqref{eq:FiniteTimeLThreeEscape}.  No such \(T<T_*\) can exist, which proves
\eqref{eq:BodySmallDataCaptureHorizon}.
\end{proof}

The kinetic-energy identity in \Cref{prop:ClassicalKineticEnergyIdentity}
supplies only \(\int_0^{T_*}Z(t)\,dt<\infty\).  Interpolating
\(D\leq Z^{1/2}Y^{1/2}\) in
\eqref{eq:BodyFirstRectangleStretchingEstimate} shows the gap directly:
\begin{equation}\label{eq:BodyCubicEnstrophyObstruction}
  Z'(t)+\nu Y(t)
  \leq C\nu^{-3}Z(t)^3.
\end{equation}
Control of \(\int Z\) does not control \(\int Z^3\) or \(\sup Z\).  The two
calculations meet at the same boundary: neither the energy row nor the critical
row bounds the quantities in
\eqref{eq:BodyFirstRectangleSignedPrefixEquivalence} for an arbitrary pair
\((u_0,\nu)\).

Energy still tells us what a failure of that bound would have to look like.
Interpolation between the energy and enstrophy heights gives
\[
  H(t)^2
  =\frac14\norm{\Lambda^{1/2}u(t)}_2^4
  \leq\frac14\norm{u(t)}_2^2 Z(t).
\]
After the kinetic-energy identity is integrated once more in time,
\begin{equation}\label{eq:BodyCriticalHeightTimeOccupancy}
  \int_0^{T_*}H(t)^2\,dt
  \leq\frac{\norm{u_0}_2^4}{8\nu}.
\end{equation}
This is not a pointwise roof.  It says that large critical height can occupy
only a short amount of time.  The interpolation between the energy height and
the critical dissipation makes the simultaneous spatial cost visible as well:
for a nonzero datum,
\begin{equation}\label{eq:BodyCriticalDissipationLowerFromHeight}
  \norm{\Lambda^{1/2}u(t)}_2
  \leq
  \norm{u(t)}_2^{2/3}
  \norm{\Lambda^{3/2}u(t)}_2^{1/3}
  \quad\Longrightarrow\quad
  D(t)\geq\frac{8H(t)^3}{\norm{u_0}_2^4}.
\end{equation}

Now take a height \(M\geq H(0)\) that the history crosses on its way to
\(2M\).  Let \(t_1\) be its first arrival at \(2M\), and let \(t_0<t_1\) be
its last visit to \(M\) before that arrival.  On
\(I_M=(t_0,t_1)\), the height stays between \(M\) and \(2M\).  Equation
\eqref{eq:BodyCriticalHeightTimeOccupancy} gives
\begin{equation}\label{eq:BodyCriticalDoublingIntervalLength}
  |I_M|
  \leq\frac{\norm{u_0}_2^4}{8\nu M^2}.
\end{equation}
At the same time,
\(|\mathcal P_{1/2}|\leq C\sqrt{2H}\,D\).  Integrating the critical balance
across this first-hitting interval gives
\begin{equation}\label{eq:BodyCriticalDoublingDissipationCost}
  M
  =\int_{I_M}H'(t)\,dt
  \leq2C\sqrt M\int_{I_M}D(t)\,dt,
  \qquad
  \int_{I_M}D(t)\,dt\geq\frac{\sqrt M}{2C}.
\end{equation}
This is the rigorous Zeno remnant.  Higher records are forced into shorter
intervals, while each completed doubling carries a larger lower amount of
critical dissipation.  The first fact leaves room in finite time; the second is
a lower cost, not the missing upper budget for \(D_*\).

\Needspace{9\baselineskip}
That unpaid budget cannot remain at a fixed range of frequencies.  For
\(K\geq1\), set
\[
  D_{\leq K}(t)
  :=\int_{|\xi|\leq K}|\xi|^3|\widehat u(t,\xi)|^2\,d\xi.
\]
Since \(|\xi|^3\leq K|\xi|^2\) on this region, energy pays
\begin{equation}\label{eq:BodyFixedFrequencyCriticalDissipation}
  \int_0^{T_*}D_{\leq K}(t)\,dt
  \leq K\int_0^{T_*}Z(t)\,dt
  \leq\frac{K\norm{u_0}_2^2}{2\nu}<\infty.
\end{equation}
If \(D_*\) diverges, its unpaid part must move beyond every fixed \(K\).
The missing estimate is an ultraviolet control on the same narrowing history,
not a failure of low frequencies or of overlap compatibility.

The size of this first gap is visible from what the first rectangle would do.
The adjacent trace lemma turns \eqref{eq:BodyFirstWholeTerminalRectangle} into
\[
  u\in C([0,T_*];H^1).
\]
Since \(H^1(\R^3)\hookrightarrow L^3(\R^3)\), it follows that
\begin{equation}\label{eq:EndpointCriterionFromLowRectangle}
  \sup_{0\leq t\leq T_*}\norm{u(t)}_3<\infty.
\end{equation}
But the endpoint theorem of Escauriaza, Seregin, and \v{S}ver\'ak gives the
opposite conclusion at a finite maximal time
\parencite{EscauriazaSereginSverak2003}:
\begin{equation}\label{eq:FiniteTimeLThreeEscape}
  T_*<\infty
  \quad\Longrightarrow\quad
  \sup_{0<t<T_*}\norm{u(t)}_3=\infty.
\end{equation}
So if the first rectangle reaches the endpoint, it already rules out the
alleged last time.  The larger intersection is not needed to repeat that
contradiction; its different work is to carry every finite response and spatial
derivative of the same history to one common terminal datum.

To make that larger conclusion, the same cutoff-independent control is needed
at every finite coordinate.  We state that remaining input at the point where
the argument begins to use it.
\begin{assumption}[Whole-terminal rectangle condition]\label[assumption]{ass:BodyWholeTerminalRectangleCondition}
Let \(u_0\in\Ssigma\), let \(\nu>0\), and let \((u,p)\) be the maximal classical
solution generated by \((u_0,\nu)\) on \([0,T_*)\).  If \(T_*<\infty\), then for
every finite \(N,M\in\Nzero\),
\begin{equation}\label{eq:BodyWholeTerminalRectangleEstimate}
  \sup_{0<T<T_*}\mathfrak R_{N,M}(u;T)<\infty.
\end{equation}
The associated pressure rows are generated by
\eqref{eq:BodyTemporalPressureRecurrence} with the fixed decaying normalization.
\end{assumption}
Here \((N,M)\) is fixed before the supremum over \(T<T_*\) is taken.  The bound
may depend on that pair; no constant uniform over all derivative orders is
asserted.  Every rectangle belongs to the same history, and larger rectangles
return the identical derivatives on their overlaps.

\Needspace{10\baselineskip}
Suppose now that the maximal history ends at a finite \(T_*\).  On this alleged
branch, \Cref{ass:BodyWholeTerminalRectangleCondition} supplies
\eqref{eq:BodyWholeTerminalRectangleEstimate} on that branch.

Fix \(N\) and \(M\).
Every term in \eqref{eq:BodyRectangleNorm} is an integral of a nonnegative squared norm.
For any one integrand \(f\), the intervals \((0,T)\) exhaust \(I_*\), so
\[
  \int_{I_*}f(t)\,dt
  =
  \lim_{T\uparrow T_*}\int_0^T f(t)\,dt.
\]
The whole-terminal estimate keeps that limit finite for every control in the
chosen rectangle:
\begin{equation}\label{eq:BodyWholeTerminalAdjacentRows}
  V_n\in L^2(I_*;H^{M+1}),
  \qquad
  V_{n+1}=\partial_tV_n\in L^2(I_*;H^{M-1}),
  \qquad 0\leq n\leq N.
\end{equation}
The pressure part of
\Cref{thm:WholeTerminalCompatibleRectangleEstimate}, using the Riesz bounds and
finite-order Sobolev products, gives the matching rows
\begin{equation}\label{eq:BodyWholeTerminalPressureRows}
  Q_n\in L^1(I_*;H^M),
  \qquad 0\leq n\leq N.
\end{equation}
Together these velocity and pressure responses form
\(\cR_{N,M}[u_0]\) on the whole interval \(I_*\).

Now enlarge the rectangle.  The fluid does not acquire a second velocity, and
a response already reached in the smaller rectangle does not change its value.
If \(N'\leq N\) and \(M'\leq M\), let
\(\rho_{N',M'}^{N,M}\) forget the extra rows and read the retained ones at the
lower Sobolev height.  Every retained entry is still the same derivative of
the same \(u\), so
\begin{equation}\label{eq:BodyRectangleRestrictionCompatibility}
  \rho_{N',M'}^{N,M}\cR_{N,M}[u_0]
  =\cR_{N',M'}[u_0].
\end{equation}
That equality is what lets the intersection be formed: every finite view of
the history agrees with every larger view wherever they see the same response.
The whole compatible family is
\begin{equation}\label{eq:DatumSpecificCompatibleIntersection}
  \cI[u_0]
  :=\bigl(\cR_{N,M}[u_0]\bigr)_{N,M\in\Nzero},
  \qquad
  \rho_{N',M'}^{N,M}\cR_{N,M}[u_0]
  =\cR_{N',M'}[u_0].
\end{equation}
The notation creates no new field.  Given any finite temporal order and
spatial height, choose one rectangle large enough to see that derivative;
every larger rectangle returns the same derivative of the same \(u\).

\begin{proposition}[The mixed intersection]\label[proposition]{prop:BodyCompatibleProjectiveAssembly}
On the finite branch just fixed, the velocity component of \(\cI[u_0]\)
satisfies
\begin{equation}\label{eq:BodyProducedMixedIntersection}
  u\in
  \bigcap_{r=0}^{\infty}
  \bigcap_{m=0}^{\infty}
  H^r(I_*;H^m(\R^3)).
\end{equation}
In particular,
\begin{equation}\label{eq:CompleteSpatialRow}
  u\in\bigcap_{m=0}^{\infty}L^2(I_*;H^m(\R^3)).
\end{equation}
\end{proposition}

\begin{proof}
Fix finite \(r,m\), and choose \(N\geq r\) and \(M\geq m\).
The upper edge of this one rectangle gives
\[
  \partial_t^nu\in L^2(I_*;H^{M+1})
  \hookrightarrow L^2(I_*;H^m)
  \qquad(0\leq n\leq r).
\]
These are the distributional time derivatives of the same velocity field, hence \(u\in H^r(I_*;H^m)\).
Because \(r\) and \(m\) were arbitrary, every finite derivative belongs to the intersection; \(r=0\) gives \eqref{eq:CompleteSpatialRow}.
\end{proof}
